% ============================================================
% 02_related_work.tex
% AHSE — Adaptive Honest Selection Ensemble
% for Multi-Horizon PV Power Forecasting
% Target journal: Sustainable Energy Technologies and Assessments (SETA)
% ============================================================

This section reviews four lines of work on which AHSE builds:
learning architectures for PV forecasting, ensemble methods for energy
time series, the clearness index as a normalised target, and the use of
weather forecasts and prediction intervals.

% ------------------------------------------------------------
\subsection{Machine learning and deep learning for PV forecasting}
\label{subsec:rw_ml_dl}

Data-driven approaches now dominate short-term PV forecasting, and
recent reviews confirm the shift away from physical and statistical
models towards machine-learning and deep-learning
pipelines~\cite{gupta2025state,kumardhaked2025exploring,yu2024ref,chandel2023techniques}.
Two families coexist within this paradigm. Gradient-boosting decision
tree (GBDT) models such as LightGBM~\cite{ke2017lightgbm} and
XGBoost~\cite{chen2016xgboost}---both rooted in Friedman's gradient
boosting framework~\cite{friedman2001greedy}---train rapidly, tolerate
missing and noisy features, and remain competitive on datasets of
moderate size. Matera~\textit{et~al.}~\cite{matera2023hourly} confirm
this by showing that ambient meteorological variables fed into a
network of shallow artificial neural networks already deliver
competitive hourly forecasts across multiple latitudes, indicating that
strong PV performance does not necessarily require deep architectures.
On the deep-learning side, recurrent networks
(LSTM~\cite{hochreiter1997lstm}, GRU~\cite{cho2014gru}) were the first
popular choice for sub-hourly lead times but lose effective memory
beyond one diurnal cycle~\cite{coya2025accuracy}, whereas
Transformer-based~\cite{vaswani2017attention} PatchTST~\cite{nie2023patchtst}
and the hierarchical MLP model N-HiTS~\cite{challu2023nhits} improve
long-horizon accuracy through input patching and multi-rate
interpolation, at the cost of longer training and weaker performance
on small datasets~\cite{suresh2025benchmarking,wang2025patchtst}.
Dimitriadis~\textit{et~al.}~\cite{dimitriadis2025deep} further show
that a deep model trained across interconnected countries must
explicitly handle distribution shift to preserve accuracy, and Luo and
Zhang~\cite{luo2022adaptive} reach a similar conclusion through a
concept-drift-triggered adaptive LSTM in which re-training fires
whenever the statistical properties of the incoming data shift.
Hybrid combinations of these families~\cite{cisse2026optimized,hou2024ref,abumohsen2024combining}
further illustrate that no single base architecture is universally
optimal. These studies share one structural limitation: a single
monolithic model is asked to serve every atmospheric regime, even
though no architecture has been shown to dominate across both stable
clear-sky and variable cloud-dominated conditions.

% ------------------------------------------------------------
\subsection{Ensemble methods for energy forecasting}
\label{subsec:rw_ensemble}

Ensembles are the natural response to regime sensitivity. Simple
averaging, inverse-error weighting and Ridge-based stacking reliably
reduce variance relative to the best single model and have been widely
applied to PV forecasting in hybrid
frameworks~\cite{cisse2026optimized,hou2024ref,perera2024regional,abumohsen2024combining},
including the Yulara plant in the Northern Territory of
Australia~\cite{cisse2026optimized} (prior work by the present
authors); however, that study's hybrid component weights remain nearly
fixed at inference, preventing real reallocation of predictive
authority when the regime shifts within minutes. The present work
addresses this limitation by introducing a validation-only selection
protocol that dynamically activates or deactivates the ensemble. Dynamic
ensemble selection (DES), initially developed for general classifier
combination~\cite{cruz2018dynamic}, addresses this by adapting weights
online, either through local-competence metrics in feature space or
through learned gating networks and reinforcement
modules~\cite{zhang2024interpretable,song2025dynamic,achouri2026dynamic,konduru2025irradiance}.
Santos~\textit{et~al.}~\cite{santos2022irradiance} apply DES directly
to solar irradiance forecasting and report statistically significant
gains over static ensembles, confirming the practical value of
regime-aware model switching. Two recurring issues persist across the
DES literature, however. First, a non-trivial share of reported gains
stems from \emph{information leakage}: the combination strategy or
meta-learner is fit on the same temporal partition used for final
evaluation, which inflates performance and undermines
reproducibility~\cite{kaufman2012leakage,kapoor2023leakage}. Second,
very few studies include an explicit \emph{guard} preventing ensemble
activation when a single model already dominates; without such a
guard, an ensemble can be selected purely because it happens to match
the test partition, a symptom of selection overfitting that only
becomes visible once the model is deployed. These two issues motivate
the validation-only, statistically guarded selection introduced in
Section~\ref{sec:methodology}.

% ------------------------------------------------------------
\subsection{Clearness index and curtailment}
\label{subsec:rw_kt}

Physics-informed features offer an alternative to learned regime
detectors~\cite{pombo2022benchmarking,mayer2021extensive}. The
clearness index \(k_t = G_\mathrm{GHI}/G_\mathrm{cs}\) normalises
measured global horizontal irradiance by the clear-sky baseline
produced by models such as Ineichen--Perez~\cite{ineichen2002broadband},
here implemented via pvlib-python~\cite{holmgren2018pvlib}. This
normalisation removes the deterministic diurnal and seasonal cycles
and isolates the stochastic component driven by cloud dynamics;
Voyant~\textit{et~al.}~\cite{voyant2025importance} show that the
choice of clear-sky model materially affects short-term forecasting
accuracy, and Lauret~\textit{et~al.}~\cite{lauret2022forecasts}
demonstrate that solar forecasts built on the clear-sky index
systematically outperform those built on raw irradiance. The clear-sky reference of a PV
\emph{power} series must describe the array rather than the horizontal
irradiance: when the measured output is limited by an inverter or by
the operator, the clearness index saturates and no longer reflects the
sky. Curtailment is common in weak grids and off-grid systems, yet
public benchmarks rarely check for it.

% ------------------------------------------------------------
\subsection{Weather forecasts and prediction intervals}
\label{subsec:rw_nwp}

Beyond a few hours, statistical forecasts based on past observations
lose skill quickly and numerical weather prediction (NWP) becomes the
main source of information~\cite{mayer2021extensive,pombo2022benchmarking}.
Using NWP in a fair evaluation requires archived forecasts that were
actually available at the issue time; analyses or forecasts from later
runs leak future information. On the uncertainty side, probabilistic
PV forecasting is increasingly
studied~\cite{xiong2025probabilistic}, and conformal prediction offers
distribution-free intervals around any point
forecaster~\cite{vovk2005algorithmic,barber2021jackknife}; adaptive
conformal inference~\cite{gibbs2021adaptive} relaxes the
exchangeability assumption and has been applied to hour-ahead PV
forecasting with Transformers~\cite{suresh2025benchmarking}. Calibrating
intervals around a \emph{selected} ensemble raises the same leakage
issue as the selection itself, since in-sample validation residuals
are optimistic.

% ------------------------------------------------------------
\subsection*{Positioning of AHSE}
\label{subsec:rw_positioning}

AHSE differs from the dynamic ensembles reviewed above in three ways.
Its selection is made per block of lead times, uses the validation
period only and accepts a combination only when a day-level bootstrap
interval of the gain excludes zero, so that it falls back to the best
single model when the evidence is weak. Archived weather forecasts are
used with an explicit publication delay. Prediction intervals are
calibrated out-of-fold with respect to the selection. Unlike earlier
studies of the Yulara plant~\cite{cisse2026optimized}, including the
first version of this work, the target is checked for curtailment and
the evaluation is reported in power, on two sites, over several seeds.

